\QHeading{第1题　分别用分数和小数表示下图的涂色部分。}

\par\noindent\makebox[\linewidth][c]{
\begin{tikzpicture}[x=1bp,y=-1bp,color=black,draw=black,text=black]
\begin{scope}\clip (56.0,119.0) rectangle (538.0,282.0);
\hdhead{92}{总复习}
\hdsection{1}{143}{分别用分数和小数表示下图的涂色部分。}
\q{96}{187}{35}{13}{①}
\q{367}{187}{35}{13}{②}
\q{643}{187}{35}{13}{③}
\path[fill=bookpink] (86.000,127.333) rectangle (110.000,209.333);
\draw[line width=.4pt] (86.000,127.333)--(86.000,209.333);
\draw[line width=.4pt] (94.000,127.333)--(94.000,209.333);
\draw[line width=.4pt] (102.000,127.333)--(102.000,209.333);
\draw[line width=.4pt] (110.000,127.333)--(110.000,209.333);
\draw[line width=.4pt] (118.000,127.333)--(118.000,209.333);
\draw[line width=.4pt] (126.000,127.333)--(126.000,209.333);
\draw[line width=.4pt] (134.000,127.333)--(134.000,209.333);
\draw[line width=.4pt] (142.000,127.333)--(142.000,209.333);
\draw[line width=.4pt] (150.000,127.333)--(150.000,209.333);
\draw[line width=.4pt] (158.000,127.333)--(158.000,209.333);
\draw[line width=.4pt] (166.000,127.333)--(166.000,209.333);
\draw[line width=.6pt] (86.000,127.333)--(166.000,127.333)--(166.000,209.333)--(86.000,209.333)--(86.000,127.333);
\path[fill=bookpink] (268.000,127.333) rectangle (284.000,207.333);
\path[fill=bookpink] (284.000,167.333) rectangle (292.000,207.333);
\draw[hdgrid] (268.000,127.333)--(268.000,207.333);
\draw[hdgrid] (276.000,127.333)--(276.000,207.333);
\draw[hdgrid] (284.000,127.333)--(284.000,207.333);
\draw[hdgrid] (292.000,127.333)--(292.000,207.333);
\draw[hdgrid] (300.000,127.333)--(300.000,207.333);
\draw[hdgrid] (308.000,127.333)--(308.000,207.333);
\draw[hdgrid] (316.000,127.333)--(316.000,207.333);
\draw[hdgrid] (324.000,127.333)--(324.000,207.333);
\draw[hdgrid] (332.000,127.333)--(332.000,207.333);
\draw[hdgrid] (340.000,127.333)--(340.000,207.333);
\draw[hdgrid] (348.000,127.333)--(348.000,207.333);
\draw[hdgrid] (268.000,127.333)--(348.000,127.333);
\draw[hdgrid] (268.000,135.333)--(348.000,135.333);
\draw[hdgrid] (268.000,143.333)--(348.000,143.333);
\draw[hdgrid] (268.000,151.333)--(348.000,151.333);
\draw[hdgrid] (268.000,159.333)--(348.000,159.333);
\draw[hdgrid] (268.000,167.333)--(348.000,167.333);
\draw[hdgrid] (268.000,175.333)--(348.000,175.333);
\draw[hdgrid] (268.000,183.333)--(348.000,183.333);
\draw[hdgrid] (268.000,191.333)--(348.000,191.333);
\draw[hdgrid] (268.000,199.333)--(348.000,199.333);
\draw[hdgrid] (268.000,207.333)--(348.000,207.333);
\draw[line width=.55pt] (492.000,168.000)--(452.000,168.000);
\path[fill=bookpink] (492.000,168.000)--(459.639,144.489)--(460.914,142.827)--(462.274,141.235)--(463.716,139.716)--(465.235,138.274)--(466.827,136.914)--(468.489,135.639)--(470.214,134.453)--(472.000,133.359)--(473.840,132.360)--(475.731,131.458)--(477.665,130.657)--(479.639,129.958)--cycle;
\draw[line width=.55pt] (492.000,168.000)--(459.639,144.489);
\draw[line width=.55pt] (492.000,168.000)--(479.639,129.958);
\path[fill=bookpink] (492.000,168.000)--(504.361,129.958)--(506.335,130.657)--(508.269,131.458)--(510.160,132.360)--(512.000,133.359)--(513.786,134.453)--(515.511,135.639)--(517.173,136.914)--(518.765,138.274)--(520.284,139.716)--(521.726,141.235)--(523.086,142.827)--(524.361,144.489)--cycle;
\draw[line width=.55pt] (492.000,168.000)--(504.361,129.958);
\draw[line width=.55pt] (492.000,168.000)--(524.361,144.489);
\draw[line width=.55pt] (492.000,168.000)--(532.000,168.000);
\path[fill=bookpink] (492.000,168.000)--(524.361,191.511)--(523.086,193.173)--(521.726,194.765)--(520.284,196.284)--(518.765,197.726)--(517.173,199.086)--(515.511,200.361)--(513.786,201.547)--(512.000,202.641)--(510.160,203.640)--(508.269,204.542)--(506.335,205.343)--(504.361,206.042)--cycle;
\draw[line width=.55pt] (492.000,168.000)--(524.361,191.511);
\draw[line width=.55pt] (492.000,168.000)--(504.361,206.042);
\path[fill=bookpink] (492.000,168.000)--(479.639,206.042)--(477.665,205.343)--(475.731,204.542)--(473.840,203.640)--(472.000,202.641)--(470.214,201.547)--(468.489,200.361)--(466.827,199.086)--(465.235,197.726)--(463.716,196.284)--(462.274,194.765)--(460.914,193.173)--(459.639,191.511)--cycle;
\draw[line width=.55pt] (492.000,168.000)--(479.639,206.042);
\draw[line width=.55pt] (492.000,168.000)--(459.639,191.511);
\draw[line width=.7pt] (492.000,168.000) circle (40bp);
\q{130}{326}{160}{12}{分数（\hspace{30pt}）}
\q{130}{369}{160}{12}{小数（\hspace{30pt}）}
\q{404}{326}{160}{12}{分数（\hspace{30pt}）}
\q{404}{369}{160}{12}{小数（\hspace{30pt}）}
\q{679}{326}{160}{12}{分数（\hspace{30pt}）}
\q{679}{369}{160}{12}{小数（\hspace{30pt}）}
\hdsection{2}{426}{用简便方法计算下面各题。}
\q{98}{474}{270}{12}{$15.74-2.98$}
\q{337}{474}{270}{12}{$17.23-3.6-6.4$}
\q{571}{474}{270}{12}{$43.63-(13.63-6.42)$}
\q{98}{643}{270}{12}{$3.8+0.7-3.8+0.7$}
\q{337}{643}{270}{12}{$9.64+12.56-2.56$}
\q{571}{643}{270}{12}{$67.4-5.6+2.6-4.4$}
\hdsection{3}{824}{解决问题。}
\q{97}{867}{750}{12}{① 小明家1月份的电费是42.48元，2月份的电费比1月份多7.85元，}
\q{130}{908}{705}{12}{比3月份多11.25元。小明家3月份的电费是多少元？}
\q{97}{1025}{750}{12}{② 一个粮仓第一天上午储存小麦3.7吨，下午储存小麦3.48吨，第二天}
\q{130}{1065}{705}{12}{比第一天少储存0.8吨。第二天储存了多少吨小麦？}

\begin{scope}[color=answerred,draw=answerred,text=answerred]\ans{128.667}{217.333}{50.667}{11}{$\frac3{10}$}
\ans{128.667}{248.667}{50.667}{11}{$0.3$}
\ans{312.000}{217.333}{50.667}{11}{$\frac{25}{100}$}
\ans{312.000}{248.667}{50.667}{11}{$0.25$}
\ans{494.000}{217.333}{50.667}{11}{$\frac4{10}$}
\ans{494.000}{248.667}{50.667}{11}{$0.4$}
\ans{65.333}{337.333}{164.667}{10}{\linespread{1}\selectfont\setlength{\jot}{1pt}$\begin{aligned}&=15.74-(3-0.02)\\&=15.74-3+0.02\\&=12.74+0.02\\&=12.76\end{aligned}$}
\ans{225.333}{337.333}{164.667}{10}{\linespread{1}\selectfont\setlength{\jot}{1pt}$\begin{aligned}&=17.23-(3.6+6.4)\\&=17.23-10\\&=7.23\end{aligned}$}
\ans{380.000}{337.333}{164.667}{10}{\linespread{1}\selectfont\setlength{\jot}{1pt}$\begin{aligned}&=43.63-13.63+6.42\\&=30+6.42\\&=36.42\end{aligned}$}
\ans{65.333}{452.667}{164.667}{10}{\linespread{1}\selectfont\setlength{\jot}{1pt}$\begin{aligned}&=3.8-3.8+0.7+0.7\\&=(3.8-3.8)+(0.7+0.7)\\&=0+1.4\\&=1.4\end{aligned}$}
\ans{225.333}{452.667}{164.667}{10}{\linespread{1}\selectfont\setlength{\jot}{1pt}$\begin{aligned}&=9.64+(12.56-2.56)\\&=9.64+10\\&=19.64\end{aligned}$}
\ans{380.000}{452.667}{164.667}{10}{\linespread{1}\selectfont\setlength{\jot}{1pt}$\begin{aligned}&=67.4+2.6-5.6-4.4\\&=(67.4+2.6)-(5.6+4.4)\\&=70-10\\&=60\end{aligned}$}
\ans{86.667}{628.000}{450.000}{10}{{\setlength{\jot}{-2pt}$\begin{aligned}[t]&\phantom{={}}42.48+7.85-11.25\\&=50.33-11.25\\&=39.08\text{（元）}\end{aligned}$}\\答：3月份电费39.08元。}
\ans{86.667}{728.000}{450.000}{10}{{\setlength{\jot}{-2pt}$\begin{aligned}[t]&\phantom{={}}3.7+3.48-0.8\\&=7.18-0.8\\&=6.38\text{（吨）}\end{aligned}$}\\答：第二天储存6.38吨。}
\end{scope}
\end{scope}\pgfresetboundingbox\path[use as bounding box] (56.0,119.0) rectangle (538.0,282.0);
\end{tikzpicture}}\par
\par\noindent




\par





\par\noindent

\QHeading{第2题　用简便方法计算下面各题。}

\CalcBlock[方法一：简便计算]{1}{\begin{aligned}[t]&\phantom{={}}\color{black}15.74-2.98\\&\color{answerred}=15.74-(3-0.02)\\&\color{answerred}=15.74-3+0.02\\&\color{answerred}=12.74+0.02\\&\color{answerred}=12.76\end{aligned}}
\par
\MethodBlock{方法二：拆成整数部分和小数部分}{}{\begin{aligned}[t]&\phantom{={}}15.74-2.98\\
&=15.74-2-0.98\\
&=13.74-0.98\\
&=12.76\end{aligned}}

\MethodBlock{方法三：按相同计数单位计算}{先看成若干个百分之一，按整数加减，最后还原成小数。}{\begin{aligned}[t]&\phantom{={}}15.74-2.98\\
&=(1574-298)\div 100\\
&=1276\div 100\\
&=12.76\end{aligned}}

\MethodBlock{方法四：补差法}{从减数先补到一个整数，再补到被减数，两段差相加。}{\begin{aligned}[t]&\phantom{={}}15.74-2.98\\
&=(3-2.98)+(15.74-3)\\
&=0.02+12.74\\
&=12.76\end{aligned}}


\CalcBlock[方法一：简便计算]{2}{\begin{aligned}[t]&\phantom{={}}\color{black}17.23-3.6-6.4\\&\color{answerred}=17.23-(3.6+6.4)\\&\color{answerred}=17.23-10\\&\color{answerred}=7.23\end{aligned}}
\MethodBlock{方法二：按运算顺序直接计算（对照）}{本法只作结果核对；本题要求简便计算，前面的简便解法仍保留。}{\begin{aligned}[t]&\phantom{={}}17.23-3.6-6.4\\
&=13.63-6.4\\
&=7.23\end{aligned}}

\CalcBlock[方法一：简便计算]{3}{\begin{aligned}[t]&\phantom{={}}\color{black}43.63-(13.63-6.42)\\&\color{answerred}=43.63-13.63+6.42\\&\color{answerred}=30+6.42\\&\color{answerred}=36.42\end{aligned}}
\MethodBlock{方法二：按运算顺序直接计算（对照）}{本法只作结果核对；本题要求简便计算，前面的简便解法仍保留。}{\begin{aligned}[t]&\phantom{={}}43.63-(13.63-6.42)\\
&=43.63-7.21\\
&=36.42\end{aligned}}

\CalcBlock[方法一：简便计算]{4}{\begin{aligned}[t]&\phantom{={}}\color{black}3.8+0.7-3.8+0.7\\&\color{answerred}=3.8-3.8+0.7+0.7\\&\color{answerred}=(3.8-3.8)+(0.7+0.7)\\&\color{answerred}=0+1.4\\&\color{answerred}=1.4\end{aligned}}
\MethodBlock{方法二：交换、结合后先凑整}{保留每个数原来的加减号，先把能凑整的部分放在一起。}{\begin{aligned}[t]&\phantom{={}}3.8+0.7-3.8+0.7\\
&=(3.8-3.8)+0.7+0.7\\
&=0+0.7+0.7\\
&=0.7+0.7\\
&=1.4\end{aligned}}

\MethodBlock{方法三：先处理括号或乘除，再提取相同因数}{用乘法分配律反过来计算。}{\begin{aligned}[t]&\phantom{={}}3.8+0.7-3.8+0.7\\
&=4.5-3.8+0.7\\
&=0.7+0.7\\
&=0.7\times (1+1)\\
&=0.7\times 2\\
&=1.4\end{aligned}}

\CalcBlock[方法一：简便计算]{5}{\begin{aligned}[t]&\phantom{={}}\color{black}9.64+12.56-2.56\\&\color{answerred}=9.64+(12.56-2.56)\\&\color{answerred}=9.64+10\\&\color{answerred}=19.64\end{aligned}}
\MethodBlock{方法二：按运算顺序直接计算（对照）}{本法只作结果核对；本题要求简便计算，前面的简便解法仍保留。}{\begin{aligned}[t]&\phantom{={}}9.64+12.56-2.56\\
&=22.2-2.56\\
&=19.64\end{aligned}}

\CalcBlock[方法一：简便计算]{6}{\begin{aligned}[t]&\phantom{={}}\color{black}67.4-5.6+2.6-4.4\\&\color{answerred}=67.4+2.6-5.6-4.4\\&\color{answerred}=(67.4+2.6)-(5.6+4.4)\\&\color{answerred}=70-10\\&\color{answerred}=60\end{aligned}}
\MethodBlock{方法二：先合并要减去的部分}{连续减去几个数，可以先把这些数相加再减。}{\begin{aligned}[t]&\phantom{={}}67.4-5.6+2.6-4.4\\
&=67.4+2.6-(5.6+4.4)\\
&=67.4+2.6-10\\
&=70-10\\
&=60\end{aligned}}

\MethodBlock{方法三：交换、结合后先凑整}{保留每个数原来的加减号，先把能凑整的部分放在一起。}{\begin{aligned}[t]&\phantom{={}}67.4-5.6+2.6-4.4\\
&=(67.4+2.6)-5.6-4.4\\
&=70-5.6-4.4\\
&=64.4-4.4\\
&=60\end{aligned}}

\MethodBlock{方法四：按运算顺序直接计算（对照）}{本法只作结果核对；本题要求简便计算，前面的简便解法仍保留。}{\begin{aligned}[t]&\phantom{={}}67.4-5.6+2.6-4.4\\
&=61.8+2.6-4.4\\
&=64.4-4.4\\
&=60\end{aligned}}

\par\noindent

\QHeading{第3题　解决问题。}

\Problem{① 小明家1月份的电费是42.48元，2月份的电费比1月份多7.85元，比3月份多11.25元。小明家3月份的电费是多少元？}
\par


\MethodTitle{方法一：简便计算}
\Answer{{\Formula{\begin{aligned}[t]&\phantom{={}}42.48+7.85-11.25\\&=42.48-(11.25-7.85)\\&=42.48-3.4\\&=39.08\text{（元）}\end{aligned}}}\par 答：3月份电费39.08元。}

\MethodTitle{方法二：直接比较3月和1月的电费}
\Answer{2月比1月多7.85元，比3月多11.25元，所以3月比1月少3.4元。}
\Answer{\Formula{\begin{aligned}[t]&\phantom{={}}42.48-(11.25-7.85)\\
&=42.48-3.4\\
&=39.08\text{（元）}\end{aligned}}}
\Answer{答：3月份电费是39.08元。}
\Problem{② 一个粮仓第一天上午储存小麦3.7吨，下午储存小麦3.48吨，第二天比第一天少储存0.8吨。第二天储存了多少吨小麦？}

\MethodTitle{方法一}
\Answer{{\Formula{\begin{aligned}[t]&\phantom{={}}3.7+3.48-0.8\\&=7.18-0.8\\&=6.38\text{（吨）}\end{aligned}}}\par 答：第二天储存6.38吨。}

\MethodTitle{方法二：先从上午量中减去差额，再加下午量}
\Answer{总和减0.8吨，可以先从上午的3.7吨里减0.8吨，再加下午量。}
\Answer{\Formula{\begin{aligned}[t]&\phantom{={}}(3.7-0.8)+3.48\\
&=2.9+3.48\\
&=6.38\text{（吨）}\end{aligned}}}
\Answer{答：第二天储存6.38吨小麦。}
